\begin{resumen}
	La enfermedad neumocócica, causada por \textit{Streptococcus pneumoniae}, es una causa importante de mortalidad global. En Cuba, tras la implementación de la vacuna heptavalente Quimi-Vio\textsuperscript{\textregistered} 7 en niños, se desarrolló el candidato vacunal Quimi-Vio\textsuperscript{\textregistered} 11 para adultos de 50 a 74 años. El presente trabajo tuvo como objetivo predecir la eficacia de Quimi-Vio\textsuperscript{\textregistered} 11 mediante un estudio de inmunopuente, extrapolando la protección a partir de la eficacia observada de Quimi-Vio\textsuperscript{\textregistered} 7. Para ello, se analizó la actividad opsonofagocítica (OPA) de 87 niños y 65 adultos un mes después de la vacunación, calculándose las Medias Geométricas de los Títulos (GMT) y su Razón (GMR) para nueve serotipos comunes. Ante valores por debajo del límite de cuantificación, se compararon el método de sustitución simple y el modelo de regresión Tobit (máxima verosimilitud, MLE). La no inferioridad del grupo adulto respecto al pediátrico se estableció cuando el límite inferior del intervalo de confianza del 95\% para la GMR fue superior a 0,5.  Los resultados demostraron no inferioridad en siete serotipos (1, 5, 6A, 19A, 6B, 18C y 19F). Los serotipos 14 y 23F no superaron el margen por una mayor respuesta pediátrica y variabilidad de los datos. El método MLE redujo el sesgo y proporcionó intervalos más robustos que la sustitución ante la censura. Se concluye que el inmunopuente permite inferir la protección de Quimi-Vio\textsuperscript{\textregistered} 11 en adultos para la mayoría de los serotipos, y que la elección del método estadístico determina la validez de las inferencias con datos censurados.
\end{resumen}

\begin{abstract}
	Pneumococcal disease, caused by \textit{Streptococcus pneumoniae}, is a major cause of global mortality. In Cuba, following the implementation of the heptavalent vaccine Quimi-Vio\textsuperscript{\textregistered} 7 in children, the vaccine candidate Quimi-Vio\textsuperscript{\textregistered} 11 was developed for adults aged 50 to 74 years. The present study aimed to predict the efficacy of Quimi-Vio\textsuperscript{\textregistered} 11 through an immunobridging study, extrapolating protection from the observed efficacy of Quimi-Vio\textsuperscript{\textregistered} 7. To this end, opsonophagocytic activity (OPA) was analyzed in 87 children and 65 adults one month after vaccination, calculating Geometric Mean Titers (GMT) and their Ratio (GMR) for nine common serotypes. For values below the limit of quantification, the simple substitution method and the Tobit regression model (maximum likelihood estimation, MLE) were compared. Non-inferiority of the adult group relative to the pediatric group was established when the lower limit of the 95\% confidence interval for the GMR was greater than 0.5. The results demonstrated non-inferiority in seven serotypes (1, 5, 6A, 19A, 6B, 18C, and 19F). Serotypes 14 and 23F did not exceed the margin due to a higher pediatric response and data variability. The MLE method reduced bias and provided more robust intervals than substitution in the presence of censoring. It is concluded that immunobridging allows inferring protection of Quimi-Vio\textsuperscript{\textregistered} 11 in adults for most serotypes, and that the choice of statistical method determines the validity of inferences with censored data.
\end{abstract}